\chapter{W5H Metadata: What}
\label{chap:W5H_What}

\section{Purpose}

The \emph{What} field defines the essence of the Engram: the event, fact,
principle, action, or conceptual payload being encoded.

\section{Requirements}

\begin{itemize}
    \item The Engram MUST provide a clear and minimally sufficient description.
    \item The \textbf{What} field must be interpretable independently of context.
    \item Ambiguity is allowed only if formally labeled (e.g., ``hypothetical'', 
          ``uncertain'', ``branch-dependent'').
\end{itemize}

\section{Schema}

\begin{quote}
\textbf{What Schema}\\
\texttt{what}: 
\begin{itemize}
    \item description: required (natural-language content)
    \item type: optional (\{fact, event, narrative, principle, structure\})
    \item scope: optional (domain or subsystem)
\end{itemize}
\end{quote}

\section{Canonicalization Rules}

\begin{itemize}
    \item \textbf{What} must be consistent with:
          \begin{itemize}
              \item ZED Core constraints,
              \item Engram type (Mote, Constructor, W5H-applied),
              \item Continuity Fields (C/O/P-class rules).
          \end{itemize}
\end{itemize}
